\chapter{Descrizione dello stage}
\label{cap:descrizione-stage}

\intro{Questa sezione è dedicata alla descrizione dell'esperienza di stage, con una visione dettagliata sulle attività previste, i rischi e i requisiti richiesti.}\\

\section{Il progetto AI Consistency Platform}
Il progetto, denominato AI Consistency Platform, consiste nello sviluppo di un'applicazione web che integra il modello di intelligenza artificiale Claude Sonnet di Anthropic per la verifica della conformità documentale rispetto a policy aziendali e requisiti normativi. \\
L'idea del progetto nasce dalla necessità di automatizzare e velocizzare il processo di verifica documentale, un processo lungo e soggetto a errori umani.
Il flusso dell'applicazione ha origine dall'interfaccia utente, attraverso la quale un utente registrato può accedere alla propria dashboard personale. \\
Dalla dashboard è possibile avviare una nuova analisi caricando un documento in formato PDF, MD o DOCX e scegliendo la tipologia del documento caricato. Sono stati implementati 5 template diversi, progettati per coprire diverse parti delle aree di documentazione aziendale. 
Il documento viene elaborato da un parser che analizza il documento ed estrae tutto il testo e le eventuali tabelle presenti. Le informazioni estratte vengono inoltrate al modulo AI, il quale genera un prompt destinato a Claude Sonnet. \\
Il modello, grazie a una fase intensa di prompt engineering, ritorna un report strutturato in cui si trovano tutte le informazioni che riguardano le sezioni del documento, incluse eventuali non confirmità rilevate.
Infine, l'utente ha la possibilità di scaricare il report in formato PDF oppure JSON, e il file originale viene salvato nel database in caso l'utente volesse riscaricarlo nel futuro. 

\section{Gli obiettivi}
Siccome il progetto è stato assegnato a due persone, ognuno dei membri aveva obiettivi diversi. I miei obiettivi sono stati:
\begin{itemize}
    \item Implementare il modulo AI per la verifica della conformità rispetto a policy interne e requisiti normativi
    \item Sviluppare il generatore automatico di report strutturati di qualità documentale
    \item Progettare e implementare la migrazione del sistema verso AWS private cloud
    \item Validare il sistema garantendo la continuità del servizio durante la transizione cloud
\end{itemize}
Mentre invece, gli obiettivi comuni erano la progettazione del sistema e dell'architettura, l'analisi dei requisiti e la creazione delle API.

\section{Il way of working}
Il progetto è stato sviluppato seguendo una modalità di lavoro collaborativa tra i due membri del team. Le attività sono state organizzate in modo iterativo, con cicli di sviluppo progressivi che includevano analisi, progettazione, implementazione, test e revisione delle funzionalità.\\
\subsection{Comunicazione e versionamento}
La comunicazione tra i membri del team è stata continua e supportata da strumenti di versionamento del codice e tracciamento delle attività. \\
Il versionamento del codice è stato gestito attraverso una \gls{monorepog} su \gls{githubg}, basata su un \gls{boilerplateg} aziendale già predisposto. 
\section{Valutazione dei rischi}
La seguente tabella indica tutti i rischi che sono stati individuati prima di iniziare la fase di sviluppo:

\begin{table}[h]
    \caption{Tabella dei rischi individuati}
\label{tab:rischi}
\begin{tabularx}{\textwidth}{XllX}
\hline\hline
\textbf{Rischio} & \textbf{Probabilità} & \textbf{Impatto} & \textbf{Mitigazione} \\
\hline
Analisi Large Language \newline Model incompleta \newline (falsi negativi)  & Alta  & Alto  & Suddivisione del \newline documento in sezioni per \newline l'analisi, anziché l'invio \newline dell' intero testo \\
\hline
Analisi Large Language \newline Model imprecisa (falsi \newline positivi) & Media & Medio & Prompt strutturati con \newline istruzioni esplicite sulla \newline soglia di segnalazione \\
\hline
Finestra di contesto \newline insufficiente per documenti di grandi dimensioni & Media & Alto & Suddivisione del documento in sezioni prima dell'invio al modello \\
\hline
Risposte del modello verbose o poco precise & Media & Medio & Limitazione della lunghezza massima della risposta per singola sezione \\
\hline
Imprecisione del quality \newline score nella valutazione \newline complessiva del documento & Media & Medio & Classificazione degli errori in tre livelli di gravità \newline (Bassa, Media e Grave) con pesi distinti nel calcolo dello score\\
\hline
Eterogeneità dei formati \newline documentali & Media & Medio & Utilizzo di librerie per ogni formato (PDF, Word, Markdown) \\
\hline
Indisponibilità o costo \newline elevato delle API Anthropic & Bassa & Alto & Limitare le chiamate alle API durante lo sviluppo \newline oppure impostare un limite massimo sui token \\
\hline
\end{tabularx}

\end{table}